\scanpage{38}
\mspage{12}
\uline{Récapitulation}

Soit \deleted{\unsure{toute}}\inserted{$X$} surface conforme (sans bord), \deleted{\unsure{on peut pas}}\inserted{On peut} définir sur $X$ une métrique riemannienne canonique $Q_X$,\inserted{compatible avec str. conforme} sauf dans les deux cas suivants:

a) $X\simeq\mathbb C$\inserted{(i.e. $X$ droite affine sur un corps\inserted{top} $C$ isom. à $\mathbb C$)}, alors on peut définir quand même une classe de métriques riemanniennes compatibles, mod. multiplication par un scalaire. Ce sont les métriques stables par translation, qui sont aussi quasi-inv. par tt autom. conforme de $X$. Elles sont \uline{plates} (donc à courbure nulle).

b) $X\simeq\mathbb P^1_{\mathbb C}$ (i.e. $X$ droite projective sur un corps top $C$ isom. à $\mathbb C$, ou encore $X$ sphère conforme). Alors il n'existe pas de métriques riemanniennes sur $X$ quasi-invariantes sous $G=\operatorname{Aut}X$, ni $=$ sous $\operatorname{Aut}^+(X)$.

\deleted{\unsure{Donc}} Il y a pourtant sur $X$ un \deleted{\unsure{tore}}\inserted{ens} $R(X)$ des métriques riemanniennes (dites ``sphériques normalisées'') avec $\int_X dx=1$, qui forment une orbite sous $G$. Si $H\subset G$ est un sous-groupe, il existe une $Q\in R(X)$ stable par $H$ ssi $H$ est rel. cpt; et $Q$ est unique ssi $X^H=\varnothing$. [Dans le cas $X^H\ne\varnothing$, cf proposition précédente qui détermine $R(X)^H$ dans tous les cas.]

\margin{\uline{NB} $R(X)$ est un espace \unsure{homogène} de dim $3$}
% archive foliation: 40
